\subsection{Clifford Gates}

\subsubsection{Definition}

In the previous section, we introduced the Heisenberg representation of quantum circuits. There, we saw that applying a Hadamard gate $H$ and then measuring an observable $Z$ is equivalent to first transforming the observable $Z$ into $X$ and then measuring $X$:
\begin{equation*}
    Z
    \;\xrightarrow{H}\;
    X
\end{equation*}
This naturally leads to an important question: \textbf{\emph{what kinds of gates transform simple observables into other simple observables?}} A particularly important class of gates with this property is known as the \textbf{Clifford gates}.

\begin{definitionbox}[: Clifford Gate]
    A unitary gate $U$ is called a \definition{Clifford Gate} if, for every Pauli operator $P$ (i.e., $P \in \{I, X, Y, Z\}$), the transformed operator
    \begin{equation}
        U^\dagger P U
        =
        P'
    \end{equation}
    Is still a Pauli operator (possibly with an overall sign $\pm 1$).

    \highspace
    In other words, \hl{Clifford gates map Pauli operators to other Pauli operators under conjugation.}
\end{definitionbox}

\noindent
For a single qubit, the Pauli operators are:
\begin{equation*}
    I\footnote{Technically, Wolfgang Pauli only defined $X$, $Y$, and $Z$ as the Pauli operators because they are the non-trivial ones. However, we include the identity $I$ in the set of Pauli operators for completeness.},
    \qquad
    X,
    \qquad
    Y,
    \qquad
    Z.
\end{equation*}
Therefore, a single-qubit Clifford gate must transform each of these operators into another operator from the same set (up to sign).

\highspace
For example:
\begin{equation*}
    Z
    \;\longrightarrow\;
    X
\end{equation*}
Is allowed, because $X$ is still a Pauli operator. However:
\begin{equation*}
    Z
    \;\longrightarrow\;
    \frac{X + Y}{\sqrt{2}}
\end{equation*}
Is not allowed, because the \textbf{result is not a single Pauli operator}.

\highspace
\begin{flushleft}
    \textcolor{Green3}{\faIcon{book} \textbf{Multi-Qubit Extension}}
\end{flushleft}
For $n$ qubits, Pauli operators are formed by taking tensor products of single-qubit Pauli matrices. For two qubits, examples of Pauli operators include:
\begin{equation*}
    X \otimes I,
    \qquad
    I \otimes Z,
    \qquad
    Y \otimes X,
    \qquad
    Z \otimes Z
\end{equation*}
For three qubits, examples include:
\begin{equation*}
    X \otimes I \otimes Z,
    \qquad
    Y \otimes Y \otimes X
\end{equation*}
In general, an $n$-qubit Pauli operator has the form:
\begin{equation*}
    P = P_1 \otimes P_2 \otimes \cdots \otimes P_n
\end{equation*}
Where each $P_i$ is one of the single-qubit Pauli operators $\left\{I, X, Y, Z\right\}$.

\highspace
\begin{examplebox}[: Hadamard Gate is a Clifford Gate]
    The Hadamard gate $H$ is:
    \begin{equation*}
        H = \frac{1}{\sqrt{2}}
        \begin{bmatrix}
            1 & 1 \\
            1 & -1
        \end{bmatrix}
    \end{equation*}
    Because $H$ is both Hermitian and unitary, we have $H^\dagger = H$ and $H^{-1} = H$. Therefore, we can compute the conjugation of the Pauli operators by $H$:
    \begin{equation*}
        H P H = P'
    \end{equation*}
    In practice, the Hadamard gate transforms the Pauli matrices as follows:
    \begin{equation}\label{eq:hadamard-conjugation}
        \begin{aligned}
            H X H &= 
            \frac{1}{2}
            \left(\begin{bmatrix}
                1 & 1 \\
                1 & -1
            \end{bmatrix}
            \begin{bmatrix}
                0 & 1 \\
                1 & 0
            \end{bmatrix}
            \begin{bmatrix}
                1 & 1 \\
                1 & -1
            \end{bmatrix}\right) \\
            &=
            \frac{1}{2}
            \left(\begin{bmatrix}
                1 & 1 \\
                -1 & 1
            \end{bmatrix}
            \begin{bmatrix}
                1 & 1 \\
                1 & -1
            \end{bmatrix}\right) \\
            &=
            \frac{1}{2}
            \left(\begin{bmatrix}
                2 & 0 \\
                0 & -2
            \end{bmatrix}\right) \\
            &=
            Z
        \end{aligned}
    \end{equation}
    \begin{equation}\label{eq:hadamard-conjugation-2}
        \begin{aligned}
            H Z H &= 
            \frac{1}{2}
            \left(\begin{bmatrix}
                1 & 1 \\
                1 & -1
            \end{bmatrix}
            \begin{bmatrix}
                1 & 0 \\
                0 & -1
            \end{bmatrix}
            \begin{bmatrix}
                1 & 1 \\
                1 & -1
            \end{bmatrix}\right) \\
            &=
            \frac{1}{2}
            \left(\begin{bmatrix}
                1 & -1 \\
                1 & 1
            \end{bmatrix}
            \begin{bmatrix}
                1 & 1 \\
                1 & -1
            \end{bmatrix}\right) \\
            &=
            \frac{1}{2}
            \left(\begin{bmatrix}
                0 & 2 \\
                2 & 0
            \end{bmatrix}\right) \\
            &=
            X
        \end{aligned}
    \end{equation}
    \begin{equation}\label{eq:hadamard-conjugation-3}
        \begin{aligned}
            H Y H &= 
            \frac{1}{2}
            \left(\begin{bmatrix}
                1 & 1 \\
                1 & -1
            \end{bmatrix}
            \begin{bmatrix}
                0 & -i \\
                i & 0
            \end{bmatrix}
            \begin{bmatrix}
                1 & 1 \\
                1 & -1
            \end{bmatrix}\right) \\
            &=
            \frac{1}{2}
            \left(\begin{bmatrix}
                i & -i \\
                -i & -i
            \end{bmatrix}
            \begin{bmatrix}
                1 & 1 \\
                1 & -1
            \end{bmatrix}\right) \\
            &=
            \frac{1}{2}
            \left(\begin{bmatrix}
                0 & 2i \\
                -2i & 0
            \end{bmatrix}\right) \\
            &=
            -Y
        \end{aligned}
    \end{equation}
    Each result is still a Pauli operator (up to a sign), confirming that the \hl{Hadamard gate is indeed a Clifford gate.}
\end{examplebox}

\highspace
\begin{examplebox}[: CNOT Gate is a Clifford Gate]
    The CNOT gate (\autopageref{sec:main-multi-qubit-gates}) is:
    \begin{equation*}
        \text{CNOT} =
        \begin{bmatrix}
            1 & 0 & 0 & 0 \\
            0 & 1 & 0 & 0 \\
            0 & 0 & 0 & 1 \\
            0 & 0 & 1 & 0
        \end{bmatrix}
    \end{equation*}
    As with the Hadamard gate, we can verify that the CNOT gate is a Clifford gate by checking how it transforms the two-qubit Pauli operators. To do this, we compute the tensor products of the single-qubit Pauli operators with the identity (to create two-qubit Pauli operators) and then conjugate them by the CNOT gate (since the CNOT gate is hermitian, we have $\text{CNOT}^\dagger = \text{CNOT}$):
    \begin{equation*}
        \begin{aligned}
            X \otimes I &= 
            \begin{bmatrix}
                0 \cdot \begin{bmatrix}
                    1 & 0 \\
                    0 & 1
                \end{bmatrix} &
                1 \cdot \begin{bmatrix}
                    1 & 0 \\
                    0 & 1
                \end{bmatrix} \\
                1 \cdot \begin{bmatrix}
                    1 & 0 \\
                    0 & 1
                \end{bmatrix} &
                0 \cdot \begin{bmatrix}
                    1 & 0 \\
                    0 & 1
                \end{bmatrix}
            \end{bmatrix} \\
            &=
            \begin{bmatrix}
                0 & 0 & 1 & 0 \\
                0 & 0 & 0 & 1 \\
                1 & 0 & 0 & 0 \\
                0 & 1 & 0 & 0
            \end{bmatrix}
        \end{aligned}
    \end{equation*}
    \begin{equation}\label{eq:cnot-conjugation}
        \begin{aligned}
            \text{CNOT} (X \otimes I) \text{CNOT} &= \begin{bmatrix}
                1 & 0 & 0 & 0 \\
                0 & 1 & 0 & 0 \\
                0 & 0 & 0 & 1 \\
                0 & 0 & 1 & 0
            \end{bmatrix}
            \begin{bmatrix}
                0 & 0 & 1 & 0 \\
                0 & 0 & 0 & 1 \\
                1 & 0 & 0 & 0 \\
                0 & 1 & 0 & 0
            \end{bmatrix}
            \begin{bmatrix}
                1 & 0 & 0 & 0 \\
                0 & 1 & 0 & 0 \\
                0 & 0 & 0 & 1 \\
                0 & 0 & 1 & 0
            \end{bmatrix} \\
            &=
            \begin{bmatrix}
                0 & 0 & 1 & 0 \\
                0 & 0 & 0 & 1 \\
                0 & 1 & 0 & 0 \\
                1 & 0 & 0 & 0
            \end{bmatrix}
            \begin{bmatrix}
                1 & 0 & 0 & 0 \\
                0 & 1 & 0 & 0 \\
                0 & 0 & 0 & 1 \\
                0 & 0 & 1 & 0
            \end{bmatrix} \\
            &=
            \begin{bmatrix}
                0 & 0 & 0 & 1 \\
                0 & 0 & 1 & 0 \\
                0 & 1 & 0 & 0 \\
                1 & 0 & 0 & 0
            \end{bmatrix} \\
            &=
            X \otimes X
        \end{aligned}
    \end{equation}
    For the $Z$ operator, the CNOT gate does not change the first qubit, so we have $Z \otimes I$:
    \begin{equation*}
        \begin{aligned}
            Z \otimes I &=
            \begin{bmatrix}
                1 \cdot \begin{bmatrix}
                    1 & 0 \\
                    0 & 1
                \end{bmatrix} &
                0 \cdot \begin{bmatrix}
                    1 & 0 \\
                    0 & 1
                \end{bmatrix} \\
                0 \cdot \begin{bmatrix}
                    1 & 0 \\
                    0 & 1
                \end{bmatrix} &
                -1 \cdot \begin{bmatrix}
                    1 & 0 \\
                    0 & 1
                \end{bmatrix}
            \end{bmatrix} \\
            &=
            \begin{bmatrix}
                1 & 0 & 0 & 0 \\
                0 & 1 & 0 & 0 \\
                0 & 0 & -1 & 0 \\
                0 & 0 & 0 & -1
            \end{bmatrix}
        \end{aligned}
    \end{equation*}
    \begin{equation}
        \begin{aligned}
            \text{CNOT} (Z \otimes I) \text{CNOT} &=
            \begin{bmatrix}
                1 & 0 & 0 & 0 \\
                0 & 1 & 0 & 0 \\
                0 & 0 & 0 & 1 \\
                0 & 0 & 1 & 0
            \end{bmatrix}
            \begin{bmatrix}
                1 & 0 & 0 & 0 \\
                0 & 1 & 0 & 0 \\
                0 & 0 & -1 & 0 \\
                0 & 0 & 0 & -1
            \end{bmatrix}
            \begin{bmatrix}
                1 & 0 & 0 & 0 \\
                0 & 1 & 0 & 0 \\
                0 & 0 & 0 & 1 \\
                0 & 0 & 1 & 0
            \end{bmatrix} \\
            &=
            \begin{bmatrix}
                1 & 0 & 0 & 0 \\
                0 & 1 & 0 & 0 \\
                0 & 0 & 0 & -1 \\
                0 & 0 & -1 & 0
            \end{bmatrix}
            \begin{bmatrix}
                1 & 0 & 0 & 0 \\
                0 & 1 & 0 & 0 \\
                0 & 0 & 0 & 1 \\
                0 & 0 & 1 & 0
            \end{bmatrix} \\
            &=
            \begin{bmatrix}
                1 & 0 & 0 & 0 \\
                0 & 1 & 0 & 0 \\
                0 & 0 & -1 & 0 \\
                0 & 0 & 0 & -1
            \end{bmatrix} \\
            &=
            Z \otimes I
        \end{aligned}
    \end{equation}
    Therefore, the CNOT gate transforms $X \otimes I$ into $X \otimes X$ and leaves $Z \otimes I$ unchanged. By performing similar calculations for all other two-qubit Pauli operators, we can confirm that the CNOT gate maps every two-qubit Pauli operator to another two-qubit Pauli operator (up to a sign):
    \begin{itemize}
        \item $\text{CNOT}\left(I \otimes X\right)\text{CNOT} = I \otimes X$
        \item $\text{CNOT}\left(I \otimes Z\right)\text{CNOT} = Z \otimes Z$
    \end{itemize}
    These results confirm that the \hl{CNOT gate is indeed a Clifford gate.}
\end{examplebox}